\section{INTRODUÇÃO}

\subsection{JUSTIFICATIVA}

***por que o estudo é relevante? Do ponto de vista prático, a relevância pode estar relacionada ao fato que as conclusões ou achados do trabalho ora realizado são úteis para a solução de problemas do dia-a-dia. Do ponto de vista teórico, um projeto de pesquisa pode se justificar em função da sua capacidade de ampliar os horizontes intelectuais da área, de propor a discussão de um tema ainda não analisado suficientemente pela ciência ou para o qual o pesquisador esteja propondo um encaminhamento distinto daquele adotado pelos seus predecessores.***

\subsection{PROBLEMA OU MOTIVAÇÃO}

***este tópico é dedicado a esclarecer o que o pesquisador pretende de fato com o seu esforço de pesquisa. Normalmente se utilizam os subitens abaixo como meios de se determinar claramente o objetivo, o que também colabora para o fechamento do escopo do trabalho***

***o problema de pesquisa, muitas vezes enunciado como uma pergunta a ser respondida, representa a questão que despertou no pesquisador a motivação para realizar a pesquisa. Está intimamente ligado ao objetivo geral, que, normalmente, consiste em encontrar a resposta para o problema de pesquisa***

\subsubsection{Motivação de Mercado}

***este tópico é dedicado ao detalhamento da motivação mercadológica do estudo do trabalho. Um guia que pode ser utilizado é o checklist do plano de negócios utilizado na monitoria das fábricas***

\subsubsection{Motivação Técnica}

***este tópico é dedicado ao detalhamento da motivação técnica do estudo do trabalho, relatando quais inovações e pesquisas estão motivando o desenvolvimento do trabalho***


\subsection{CONTRIBUIÇÕES}

***listar a contribuições geradas pelo trabalho, tais como resultados a serem obtidos, artigos publicados, entre outros.***

\subsection{OBJETIVOS}

\subsubsection{Objetivo geral}

***é o que se pretende fazer, ou seja, aquilo que, depois de atingido, fará com que o pesquisador considere que concluiu o seu trabalho***

\subsubsection{Objetivos específicos}

***são objetivos auxiliares, todos eles necessários (individualmente) e suficientes (em conjunto) para a consecução do objetivo geral. Isto quer dizer que não deve haver nenhum objetivo específico que se proponha a fazer algo que vá além do escopo do objetivo geral (critério de necessidade do objetivo específico) e que, conquistados todos os objetivos específicos, o objetivo geral é, naturalmente, também alcançado (critério de suficiência dos objetivos específicos)***

\subsection{ESTRUTURA DA DISSERTAÇÃO}

***neste item você vai descrever como está constituída a dissertação, indicando o que será encontrado em cada um dos capítulos seguintes.***